\documentclass{article}
\usepackage[utf8]{inputenc}
\usepackage[italian]{babel}
\usepackage{xcolor}
\usepackage[ruled,lined,linesnumbered]{algorithm2e}

\title{Algoritmi e strutture dati}
\author{Riccardo Di Michele}

\begin{document}
\maketitle

\section*{Introduzione}
L'obbiettivo della materia è quello di studiare la 
\textbf{complessità computazionale} degli algoritmi nel tempo; in questo 
testo non andremo a spiegare le nozioni primitive dell'informatica (esempio:
cosa è un algoritmo, una condizione if ecc.).
Un esempio di complessità computazionale:

\vspace{0.3cm}

\begin{algorithm}[H]
\SetKwData{MatA}{A}
\SetKwData{MatB}{B}
\SetKwData{MatC}{C}
\SetKwInOut{Input}{Input}
\SetKwInOut{Output}{Output}

\TitleOfAlgo{CalcoloMatriceDistanzeManhattan(\MatA, \MatB, $n$)}

\Input{Due array/vettori \MatA e \MatB ciascuno di $n$ numeri reali}
\Output{Una matrice \MatC di dimensione $n \times n$ con le differenze}

\BlankLine
\For{$i \leftarrow 1$ \KwTo $n$}{
    \For{$j \leftarrow 1$ \KwTo $n$}{
        \MatC$[i, j] \leftarrow \MatA[i] \times \MatB[j] + i \times j$\;
    }
}
\end{algorithm}

\vspace{0.3cm}
L'algoritmo appena mostrato ha costo computazionale $T(n) = n^2$, ma cosa 
vuol dire e soprattutto come usiamo
questa informazione a nostro vantaggio? \\
Sapere il costo di un algortimo ci va in nostro aiuto per vari motivi: se 
l'algoritmo è troppo costoso potremmo apportare dei cambiamenti, o addirittura
scegliere un algortimo già esistente più efficiente di quello scelto prima, o 
ancora, per capire quale è l'input più grande che l'algoritmo può "reggere" prima
di perdere troppo tempo.
\\
\\
Cominciamo allora a rispondere alla domanda più importante. \\
Definiamo la \textit{complessità computazionale} come una funzione che 
mappa la dimensione dell'input di un algortimo al suo "tempo" di calcolo
\[ costo: input \to tempo \]

Adesso dobbiamo però definire questo \textit{input} e questo \textit{tempo}. \\
Il primo lo definiamo usando il \textbf{criterio del costo uniforme}, ovvero
la dimensione dell'input dipende solo da quanti elementi ha quell'input;
se prendiamo come esempio un array di 20 elementi il suo costo sarà di 
$n = 20$. \\
Definiamo il tempo invece come il numero di operazioni elementari, ovvero 
operazioni che la macchina riesce a fare senza l'aiuto di altri operazioni.

\begin{itemize}
  \item $a = a + 1$ \textbf{operazione elementare}
  \item $a = a - 1$ \textbf{operazione elementare}
  \item $swap(A[1], A[2])$ \textbf{operazione elementare}
  \item $sin(a)$ \textbf{operazione} considerata al giorno d'oggi \textbf{elementare}
  \item $max(array, a)$ \textbf{operazione non elementare}
\end{itemize}

L'ultima operazione, ovvero la funzione $max()$ non viene considerata elementare
poichè all'interno di essa (pure se non detto in maniera esplicita) avremo 
\textbf{sicuramente} più di un'operazione elementare, come per esempio multeplici
assegnazioni di variabili.
\\ Con le informazioni acquisite possiamo adesso verificare il costo dell'algoritmo 
citato prima:

\vspace{0.3cm}

\begin{algorithm}[H]
\SetKwData{MatA}{A}
\SetKwData{MatB}{B}
\SetKwData{MatC}{C}
\SetKwInOut{Input}{Input}
\SetKwInOut{Output}{Output}

\caption{CalcoloMatriceDistanzeManhattan(\protect\MatA, \protect\MatB, $n$)}

\Input{Due array/vettori \MatA e \MatB ciascuno di $n$ numeri reali}
\Output{Una matrice \MatC di dimensione $n \times n$ con le differenze}

\BlankLine
\For{$i \leftarrow 1$ \KwTo $n$}{
    \For{$j \leftarrow 1$ \KwTo $n$}{
        \MatC$[i, j] \leftarrow \MatA[i] \times \MatB[j] + i \times j$\;
    }
}
\end{algorithm}

\vspace{0.3cm}

Per analizzare il codice cominciamo a partire dalle operazioni più interne,
prendiamo l'operazione alla riga 3

\[ C[i,j] \leftarrow A[i] \times B[j] + (i \times j) \]

Il costo di questa riga sarebbe pari al numero di operazioni elementari al 
suo interno, quindi 4; in questo caso diremmo che il \textbf{costo esatto} 
è pari a $R3 = 4$. \\
Andando a ritroso quindi analizziamo ora la parte di codice meno interna: 
il \textit{ciclo for} che va da $R2$ a $R4$. Questo è un ciclo che va da
$j \leftarrow 1$ a $j \leftarrow n$, quindi farà sempre $n$ cicli.
Col costo della riga $R3$ possiamo quindi calcolare il costo di questo ciclo come:

\[ \sum_{j \leftarrow 1}^{n} R3 \rightarrow n \cdot R3 \]

Con la stessa logica analizziamo infine il ciclo più esterno, quello tra 
$R1$ e $R5$; in questo caso la varaibile contatore è $i$, ma la logica, e il 
numero di cicli, non cambia affatto.

\[ 
\sum_{i \leftarrow 1}^{n} \sum_{j \leftarrow 1}^{n} R3 
\rightarrow n \cdot n \cdot R3 \rightarrow
n^2 \cdot R3
\]

Ma nella maggior parte dei nostri esercizi il costo esatto non ci 
interessa proprio, ci interessa invece il \textbf{costo asintotico}, ovvero
il costo dell'algoritmo considerando un input $n$ sempre più grande. \\
Poichè il costo di $R3$ è un costo \textit{fisso}, non dipendente dalla grandezza
dell'input $n$, nel \textbf{calcolo asintotico} diventa un dato insignificante,
è per questo che in tanti esercizi lo troviamo scritto come $c$; il calcolo 
di prima diventer allora:

\[ 
\sum_{i \leftarrow 1}^{n} \sum_{j \leftarrow 1}^{n} c 
\rightarrow n \cdot n \cdot c \rightarrow
c \cdot n^2 \rightarrow n^2 = T(n)
\]

Per lo stesso motivo viene direttamente omesso il valore $c$ quando 
raggiungiamo il risultato finale.


\end{document}
